% Guideline item 3: Introduction.
\section{Introduction}\label{sec:introduction}

An anomaly-alerting system has two parts. A \emph{scorer} turns each window of a
sensor signal into a number, the anomaly score; a \emph{policy} turns that score into
a decision. The simplest policy raises an alert whenever the score exceeds a fixed
threshold $\theta_0$ set on a calibration period. On a signal whose normal level drifts,
a fixed threshold goes out of date, so a common remedy is to let the threshold learn
from recent scores. This report studies one question inside that remedy:

\begin{quote}
\emph{When should an alert threshold learn, and from which data, if that data may be
chosen neither by the threshold's own decisions nor unconditionally?}
\end{quote}

We keep the scorer fixed and study the policy only. At every window end our policy
chooses one of four actions: \textsc{alert}, \textsc{hold} (no alert),
\textsc{defer} (withhold a decision, which lowers decision coverage) or
\textsc{recalibrate} (move the threshold). Recalibration is guarded: it uses a buffer of
recent scores that is filled without looking at the threshold's decisions, and it may
only move the threshold within $[\theta_0, 4\theta_0]$. All policies replay the
identical score trace, so any difference between them is caused by the decision rule
and not by a different scorer.

This is a Phase-I report. Its evidence is a development comparison on two SKAB
streams~\citep{skab2020} and a deterministic synthetic family, not a held-out result.
We therefore claim a proved mechanism and honestly scoped measurements, not that any
policy wins. Section~\ref{sec:motivation} motivates the question with a worked example;
Section~\ref{sec:literature} positions it against the closest work;
Sections~\ref{sec:problem}--\ref{sec:platform} define the problem, method and platform;
Section~\ref{sec:results} tests the claims; Sections~\ref{sec:limitations}
and~\ref{sec:conclusion} give the strongest case against us and the next experiment.
